% p090_1; source PDF page 90; manually verified editable table.
\begin{tikzpicture}[x=1cm,y=-1cm,inner sep=0pt,outer sep=0pt]
\path[use as bounding box] (0,0) rectangle (15.700,3.875);
\draw[line width=.4pt] (0.000,0.000) rectangle (3.400,0.855);
\node[align=center,text width=3.140cm,font={\linespread{1}\fontsize{10}{13.50}\selectfont}] at (1.700,0.427) {客观要件};
\draw[line width=.4pt] (3.400,0.000) rectangle (9.200,0.855);
\node[align=center,text width=5.540cm,font={\linespread{1}\fontsize{10}{13.50}\selectfont}] at (6.300,0.427) {（作为或不作为的）不法侵害};
\draw[line width=.4pt] (0.000,0.855) rectangle (3.400,2.365);
\node[align=center,text width=3.140cm,font={\linespread{1}\fontsize{10}{13.50}\selectfont}] at (1.700,1.610) {主观要件};
\draw[line width=.4pt] (3.400,0.855) rectangle (9.200,1.610);
\node[align=center,text width=5.540cm,font={\linespread{1}\fontsize{10}{13.50}\selectfont}] at (6.300,1.232) {故意的不法侵害};
\draw[line width=.4pt] (3.400,1.610) rectangle (9.200,2.365);
\node[align=center,text width=5.540cm,font={\linespread{1}\fontsize{10}{13.50}\selectfont}] at (6.300,1.987) {过失的不法侵害};
\draw[line width=.4pt] (0.000,2.365) rectangle (3.400,3.875);
\node[align=center,text width=3.140cm,font={\linespread{1}\fontsize{10}{13.50}\selectfont}] at (1.700,3.120) {排除责任事由};
\draw[line width=.4pt] (3.400,2.365) rectangle (9.200,3.120);
\node[align=center,text width=5.540cm,font={\linespread{1}\fontsize{10}{13.50}\selectfont}] at (6.300,2.742) {无责任年龄的不法侵害};
\draw[line width=.4pt] (3.400,3.120) rectangle (9.200,3.875);
\node[align=center,text width=5.540cm,font={\linespread{1}\fontsize{10}{13.50}\selectfont}] at (6.300,3.497) {无责任能力的不法侵害};
\draw[decorate,decoration={brace,amplitude=5pt},line width=.4pt] (9.65,.2) -- (9.65,3.65);
\node[anchor=west,align=left,text width=5.5cm,font=\fontsize{10}{13.5}\selectfont] at (10.05,1.925) {对其均可正当防卫};
\end{tikzpicture}
